\section*{The Propers: Detailed Commentary}

Each proper is given its appointed Latin exactly as the controlling facsimile
prints it, then an identified public-domain English translation, then the
evidence. The English throughout is Dom Laurence Shepherd's, from
\work{The Liturgical Year} \dots\ \work{Time after Pentecost}, vol.~II, 2nd
ed.\ (Stanbrook Abbey: Burns \& Oates, 1909), at the pages named in each
block. The project has translated nothing: where no identified translation is
quoted, the English is analysis and is written as analysis.

\subsection*{The text: what two 1962 books actually print}

The Latin below was collated at 300 and 600 dpi against the CMAA facsimile of
the Vatican \latin{editio typica} (printed pp.~388--389), and independently
against the page images of the Benziger \latin{editio iuxta typicam} of the
same year (printed pp.~382--383). The two books are not two photographs of
one page; they are two editions, and they differ. Nothing below is blended,
and the typical edition governs.

\begin{receptiontable}{Place}{Vatican \latin{editio typica}}{Benziger \latin{iuxta typicam}}
1 Cor. 10:8 & \latin{vigínti tria mília} & \latin{vigínti tria míllia}\\
Lk. 19:45 & \latin{cœpit eícere} & \latin{cœpit eiícere}\\
Introit close & \latin{℣. Glória Patri.} & \latin{℣. Glória Patri. Ecce Deus.}\\
Secret conclusion & \latin{Per Dóminum nostrum Iesum Christum, Fílium tuum: Qui tecum vivit et regnat in unitáte.} & \latin{Per Dóminum.}\\
Collect, Postcomm. conclusion & \latin{Per Dóminum nostrum.} & \latin{Per Dóminum.}\\
Communion reference & \latin{Io. 6, 57} & \latin{Ioann. 6, 57}\\
\end{receptiontable}

Two of these are worth more than a footnote. The Clementine Vulgate reads
\latin{míllia} at 1 Corinthians 10:8, and so does Benziger; the Vatican
typical edition prints \latin{mília}. And at Luke 19:45 the typical edition
prints the contracted \latin{eícere} where Benziger, the Clementine Vulgate
and Shepherd's Latin all print the uncontracted form. Neither changes a
meaning. Both show that ``the 1962 Latin'' is not a single artefact, and that
a guide which does not name its witness has not told the reader what it
checked.

\subsection*{Introit \cue{Int.}}

\begin{appointedlatin}{\latin{Ant. ad Introitum} --- Ps. 53, 6-7 --- marginal no.\ 1522}
Ecce Deus ádiuvat me, et Dóminus suscéptor est ánimæ meæ: avérte mala inimícis
meis, et in veritáte tua dispérde illos, protéctor meus, Dómine. \latin{Ps.
ibid., 3} Deus, in nómine tuo salvum me fac: et in virtúte tua líbera me.
℣.~Glória Patri.
\end{appointedlatin}

\begin{namedtranslation}{English: Shepherd, \work{The Liturgical Year}, Time after Pentecost II, p.~224}
Behold! God is my helper, and the Lord is the support of my soul: turn out the
evils upon mine enemies, and cut them off in thy truth, O Lord, my protector.
\textit{Ps.} O God, in thy name save me: and, in thy strength, deliver me.
Glory, \textit{etc}.
\end{namedtranslation}

\textbf{The chant text is not the Bible's text.} Set the appointed antiphon
beside the Clementine Vulgate psalter and three differences appear at once.
The Vulgate has \latin{Ecce enim Deus adjuvat me}; the Introit drops
\latin{enim}. The Vulgate's v.~3 ends \latin{et in virtute tua judica me};
the Introit sings \latin{líbera me}. And the antiphon's closing vocative
\latin{protéctor meus, Dómine} has no counterpart anywhere in Psalm 53:6--9
as the Clementine prints it.

This is not a defect in the missal. The Roman Church used an older Latin
psalter before Jerome's Gallican recension displaced it in the Bible, and the
1910 \work{Catholic Encyclopedia} records that the Roman recension ``remained
in use in Rome for a considerable time, and at the present day is still used
in the Divine Office chanted at St.~Peter's.'' That is the natural
explanation of a chant text older than the Bible text beside it. This guide
records the divergence as a verified textual observation and does not assert
the derivation, because no \work{Psalterium Romanum} witness was collated for
it.

\textbf{The consequence for reception, and it is a real one.} The Fathers and
the scholastics who expound this psalm are expounding a different verb. In
Augustine's Latin lemma the second half of v.~7 reads \latin{In virtute tua
disperde illos}, and he expounds \latin{virtus}; the Introit sings
\latin{in veritáte tua}. At v.~3 Augustine, Aquinas and Bellarmine all read
and gloss \latin{iúdica me}, not the Introit's \latin{líbera me}. A guide
that quotes their exegesis of the Introit's psalm verse without saying so has
quietly attributed to them a comment on words they did not have.

\begin{witnessnote}
\wlocus{St.\ Augustine, \work{Enarrationes in Psalmos} 53, §§4, 8--9}
The controlling image is the Ziphites as those who \textit{flourish}: their
bloom is outward and temporary, while the just man's flowering is hidden. At
§8 --- the Introit's opening clause --- Augustine's point is that the help is
\textit{interior and present}, not merely future: the glory of the just is the
testimony of conscience and the light of God's countenance ``put into my
heart --- not into my vineyard, not into my flock, not into my cask, not into
my table.'' At §9 the imprecation is read eschatologically: the enemies
flourish for a time and perish for ever. English quoted from the
public-domain NPNF first series, vol.~8, where the exposition is headed
``Psalm 54'' but cites internally as Psalm 53.
\end{witnessnote}

\begin{witnessnote}
\wlocus{St.\ Thomas Aquinas, \work{Super Psalmo} 53, nn.~1, 3--4}
Aquinas anchors the title in 1 Kings (1 Samuel) 23 and allegorises the
Ziphites in the same direction as Augustine. His structural point is that the
psalm's petition is heard twice over: against the enemies' assault,
\latin{ecce Deus adiuvat me}; against those seeking his life, \latin{Dominus
susceptor est animae meae}. He then meets the obvious objection to praying
\latin{dispérde illos} at all --- \latin{orate pro persequentibus} --- and
answers that prophetic imprecations admit three readings: as prediction (he
notes that the Hebrew has a present, ``thou turnest''), as conformity to
divine justice, and as spiritual denunciation, since ``when sinners cease to
sin, then they die and cease to be sinners.''
\end{witnessnote}

\begin{witnessnote}
\wlocus{St.\ Robert Bellarmine, \work{Explanatio in Psalmos}, Ps.~LIII, vv.~4--5, tr.\ John O'Sullivan (Dublin, 1866)}
Bellarmine turns the phrase inside out with one sentence: the enemies ``had
not God before their eyes, but God had them before his eyes.'' He reads
\latin{Ecce} as ``a sudden light from God'' and the present tense as
certainty, then supplies the concrete deliverance --- the messenger announcing
the Philistine raid that recalled Saul --- and suggests David may have been
speaking the words at that very moment. On the imprecation he agrees with
Aquinas: ``Such imprecations \dots\ are to be read as predictions; and so this
reads in the Hebrew.''
\end{witnessnote}

\begin{witnessnote}
\wlocus{Theodoret of Cyrus, \work{Interpretatio in Psalmos} 53 (PG 80, cols.\ 1265--1268) --- Greek-tradition sample}
Theodoret sets the psalm historically among the Ziphites who ``feigned
friendship'' and betrayed David's hiding place, then makes it deliberately
general: David offered it as instruction so that anyone under unjust attack
might pray it. He reads \latin{Ecce Deus ádiuvat me} as confidence in help
coming before long. Alone among the four witnesses here he reads \latin{in
veritáte tua} --- the Introit's own text --- and takes it, strikingly, as
addressed to the Holy Spirit. \textit{Source limit:} Theodoret was consulted
in Robert Hill's modern English translation, which is in copyright; he is
therefore summarised and not quoted, and the PG column numbers above are those
printed inline in that translation, not read off a Migne scan.
\end{witnessnote}

\subsection*{Collect \cue{Coll.}}

\begin{appointedlatin}{\latin{Oratio} --- marginal no.\ 1523}
Páteant aures misericórdiæ tuæ, Dómine, précibus supplicántium: et, ut
peténtibus desideráta concédas; fac eos, quæ tibi sunt plácita, postuláre. Per
Dóminum nostrum.
\end{appointedlatin}

\begin{namedtranslation}{English: Shepherd, p.~224}
May the ears of thy mercy, O Lord, be opened to the prayers of thy suppliants:
and, that thou mayst grant to thy petitioners the things they desire, make
them to ask those that are agreeable to thee. Through, \textit{etc}.
\end{namedtranslation}

The prayer is built on a single inversion. An ordinary collect asks that
petitions be heard and granted. This one asks that they be heard, and then ---
\textit{in order that} they may be granted --- asks that they first be
changed. The purpose clause \latin{ut peténtibus desideráta concédas} depends
on the imperative that follows it, \latin{fac \dots\ postuláre}, and the
printed semicolon marks the join. God is asked to do something to the
petitioners before he does anything about the petition.

Two features of the Latin are worth holding. \latin{Páteant} is not ``open
thine ears'' but ``let thine ears lie open'': a state, not an act, and the
metaphor is of a door already ajar. And \latin{desideráta} --- ``the things
desired'' --- is left deliberately unqualified, so that the prayer never says
the things desired are good. It says only that they are desired, and then
asks that what is desired be made to coincide with what pleases God.

That is why this Collect belongs in front of this Epistle. The wilderness
generation of 1 Corinthians 10 did not lack petitions; Numbers 11 is a
narrative of granted requests. \latin{Non simus concupiscéntes malórum} is
the Epistle's way of saying what the Collect asks God to prevent.

\textit{Source boundary.} No patristic or medieval commentary on this
oration's wording was located; composed Roman collects of this type carry no
ancient commentary tradition comparable to that on the Scriptures. The
analysis above rests on the printed Latin, its punctuation, and its
juxtaposition, and is offered as source-grounded synthesis rather than
documented reception.

\subsection*{Epistle \cue{Ep.}}

\begin{appointedlatin}{\latin{Léctio Epístolæ beáti Pauli Apóstoli ad Corínthios} --- 1 Cor. 10, 6-13 --- marginal no.\ 1524}
Fratres: Non simus concupiscéntes malórum, sicut et illi concupiérunt. Neque
idolólatræ efficiámini, sicut quidam ex ipsis: quemádmodum scriptum est: Sedit
pópulus manducáre et bíbere, et surrexérunt lúdere. Neque fornicémur, sicut
quidam ex ipsis fornicáti sunt, et cecidérunt una die vigínti tria mília.
Neque tentémus Christum, sicut quidam eórum tentavérunt, et a serpéntibus
periérunt. Neque murmuravéritis, sicut quidam eórum murmuravérunt, et
periérunt ab exterminatóre. Hæc autem ómnia in figúra contingébant illis:
scripta sunt autem ad correptiónem nostram, in quos fines sæculórum
devenérunt. Itaque qui se exístimat stare, vídeat ne cadat. Tentátio vos non
apprehéndat, nisi humána: fidélis autem Deus est, qui non patiétur vos tentári
supra id quod potéstis, sed fáciet étiam cum tentatióne provéntum, ut possítis
sustinére.
\end{appointedlatin}

\begin{namedtranslation}{English: Shepherd, pp.~225--226}
Brethren: Let us not covet evil things, as they also coveted. Neither become
ye idolaters, as some of them: as it is written: The people sat down to eat and
drink, and rose up to play. Neither let us commit fornication, as some of them
committed fornication, and there fell in one day three-and-twenty thousand.
Neither let us tempt Christ: as some of them tempted, and perished by the
serpents. Neither do ye murmur: as some of them murmured, and were destroyed by
the destroyer. Now, all these things happened to them in figure: and they are
written for our correction, upon whom the ends of the world are come. Wherefore
he that thinketh himself to stand, let him take heed lest he fall. Let no
temptation take hold on you, but such as is human. And God is faithful, who will
not suffer you to be tempted above that which ye are able; but will make also
with temptation issue, that ye may be able to bear it.
\end{namedtranslation}

\textbf{Where the lesson begins and where it stops.} The liturgical
\latin{Fratres:} is the missal's, not Paul's. The pericope opens mid-argument
at v.~6 --- so the assembly does not hear vv.~1--5, where the sea and the cloud
are read as a baptism and the rock is Christ --- and it closes at v.~13,
so the assembly does not hear v.~14, where Paul's \latin{propter quod} leads
straight from divine faithfulness to \latin{fugite ab idolorum cultura} and
then, at vv.~16--17, to the cup, the bread, and the many who are one body. The
Sunday will get to the one bread; it will get there through the Communion and
the Postcommunion rather than through the lesson.

\textbf{The five ruins and their sources.} Each clause of vv.~6--10 points at
a specific narrative, and one of them is quoted verbatim.

\begin{receptiontable}{Clause in the Epistle}{Narrative in view}{What the Epistle does with it}
\latin{concupiscéntes malórum} (v.~6) & Numbers 11: the craving for meat, and the graves of craving & Names desire itself, not its object, as the first ruin --- which is what the Collect has just asked God to re-form\\
\latin{Sedit pópulus manducáre et bíbere, et surrexérunt lúdere} (v.~7) & Exodus 32:6, the golden calf & The one direct citation. In the Clementine Vulgate the two run word for word --- \latin{sedit populus manducare, et bibere, et surrexerunt ludere} --- and the agreement is a fact about that text, not about every Latin recension. Note Exodus 32:5: Aaron had proclaimed \latin{cras sollemnitas Domini est}, so the meal is a cultic meal offered to the true God under an image\\
\latin{cecidérunt \dots\ vigínti tria mília} (v.~8) & Numbers 25: Baal-Peor & Numbers 25:9 reads \latin{viginti quattuor milia}, and so do the Hebrew and the Greek. No witness to Numbers gives twenty-three thousand: the difference is Paul's, and it is left standing\\
\latin{a serpéntibus periérunt} (v.~9) & Numbers 21:5--6 & The complaint there is disgust at the manna --- \latin{anima nostra iam nausiat super cibo isto levissimo} --- which is why the clause can be chained to a Eucharistic warning at all\\
\latin{periérunt ab exterminatóre} (v.~10) & The murmuring cycles & Speech is added to appetite and act. The referent is genuinely unsettled: \latin{exterminator} has no lexical counterpart in the Vulgate of either chapter, and the two Latin commentators anchor it differently --- Ambrosiaster at Numbers 14, Peter of Tarantaise at Numbers 16:41\\
\end{receptiontable}

\textbf{Twenty-three thousand.} The tradition did notice. The commentary
printed in the Aquinas corpus at this verse gives two options and refuses to
invent a third: \latin{sed maior numerus non excludit minorem, unde non
dicitur hic cum praecisione, vel forte vitium scriptorum est} --- the greater
number does not exclude the lesser, so Paul is not speaking with precision;
or perhaps it is a copyist's fault. Chrysostom and Ambrosiaster, who both
comment on the verse, do not raise the question at all. This guide reports the
difference and adopts no harmonisation.

\begin{witnessnote}
\wlocus{St.\ John Chrysostom, \work{Homilies on First Corinthians} 23 and 24}
The pericope straddles two homilies: Homily 23 runs to v.~12 and stops;
v.~13 opens Homily 24. In 23 he reads the Exodus material in strict parallel
--- as the gifts were figures, so are the punishments, and therefore ``if the
one be type, and the other substance, it must needs be that the punishments
should as far exceed as the gifts.'' He traces the whole catalogue to one
root, luxury: from sensuality they ``passed into idolatry.'' On v.~12 the sin
he identifies is not standing but \textit{thinking} one stands: ``this is not
even standing as one ought to stand, to rely on yourself \dots\ For our
standing here is not secure standing, no not until we be delivered out of the
waves of this present life.'' In 24 he reads v.~13 as a deliberate descent
after terror, and then makes the decisive move: the ability is not natural.
``For the ability lies in God's gracious influence; a power which we draw down
by our own will.'' The way through is not primarily removal: ``He gives
patience and brings on a speedy release; so that in this way also the
temptation becomes bearable.'' English from the public-domain NPNF first
series, vol.~12, tr.\ Talbot W.~Chambers.
\end{witnessnote}

\begin{witnessnote}
\wlocus{Ambrosiaster, \work{Commentaria in Epistolam ad Corinthios Primam}, ad 10:6--13 (PL 17)}
Ambrosiaster identifies the target with unusual precision: those \latin{qui in
idolio recumbentes putabant non esse peccatum} --- who reclined at table in an
idol's temple thinking it no sin. He reads v.~12 socially rather than merely
morally: it is said to those presuming on the knowledge that all food is
lawful, who scandalised weaker brethren while judging the Apostle. And at
v.~13 he gives the pericope's most distinctive reading: \latin{tentatio
humana} is not a concession but an exhortation --- \latin{Propter Christum
ergo pati, humana tentatio est, per quam proficitur apud Deum}. On the
\latin{provéntum} he is concrete to the point of arithmetic: God will either
give strength or cut the trial short, ``so that he who is known to be able to
endure no more than three days is not allowed to suffer on the fourth.''
\textit{Edition note:} read in the Migne text of the 1845 printing, where the
embedded column marks put the passage at coll.~144--145; other Migne printings
of the same work carry a different column system, and CSEL 81/2 was not
consulted.
\end{witnessnote}

\begin{witnessnote}
\wlocus{Not Aquinas: the commentary on this chapter is Peter of Tarantaise}
The exposition of 1 Corinthians 10 that circulates inside the Thomistic corpus
is not St.\ Thomas's. Corpus Thomisticum's editorial note states \latin{A
capite vero X amissa est lectura, ac lacuna repleta ex commentario Petri de
Tarantasia breviato forte a Nicholao de Gorran}, and Fabian Larcher's English
translation prints the bracket ``[CHAPTERS 7:15---10:33 (nos.\ 347--581)
supplied by Peter of Tarantaise]'' in place of the text. Any guide that
attributes this material to Aquinas at a Marietti number is attributing to him
a text his editors flag as someone else's. What the passage actually contains
is worth having on its own terms: at v.~12 it draws four things out of the
verse --- the multitude of those falling, the uncertainty of those standing,
the necessity of caution, the ease of ruin --- and at v.~13 gives
\latin{proventum} a moral condition: those do not burst in the furnace
\latin{qui non habent ventum superbiae}. One caution travels with it. Its
maxim \latin{Augustinus: propter Christum pati humana tentatio est} is
verbatim Ambrosiaster, quoted just above; it is a gloss-transmission
misattribution and should not be repeated as Augustine.
\end{witnessnote}

\textit{Negative result on Augustine.} Two loci were verified in which
Augustine uses v.~13, and both are doctrinal reuse rather than exegesis of the
pericope: \work{De sancta virginitate} 47, where untested strength is made a
ground of humility --- the virgin cannot know whether a married woman is
already able to suffer what she is merely spared, since \latin{infirmitas eius
temptatione non interrogatur} --- and \work{De bono viduitatis} 17, where the
verse is turned against those who make prayer superfluous by claiming the will
suffices. Note that his lemma reads \latin{exitum}, not the Vulgate's
\latin{provéntum}. No Augustinian citation of v.~12 was established; the
\work{Enarrationes} and \work{Sermones} were not exhaustively searched for it,
so this is ``not established,'' not ``absent.''

\textbf{The verse the twentieth century mislearned.} Verse 13 says four
things, and popular usage keeps only a garbled third. It says that the
testing which has taken hold of the Corinthians is \latin{humána} --- on
human scale, the common lot. It says the ground of confidence is God's
fidelity, not the hearer's capacity: \latin{fidélis autem Deus est}. It says
God will not permit a testing \latin{supra id quod potéstis} --- and
\latin{potéstis} is plural throughout, addressed to a community, not to an
isolated sufferer. And it says that God \textit{makes} something with the
testing, a \latin{provéntum}: an outcome, an issue, a way through, which is
supplied rather than found. Shepherd renders it ``issue.'' What the verse does
not say is that every affliction has been measured to what one person can
privately bear.

\subsection*{Gradual \cue{Grad.}}

\begin{appointedlatin}{\latin{Graduale} --- Ps. 8, 2 --- marginal no.\ 1525}
Dómine Dóminus noster, quam admirábile est nomen tuum in univérsa terra! ℣.
Quóniam eleváta est magnificéntia tua super cælos.
\end{appointedlatin}

\begin{namedtranslation}{English: Shepherd, p.~229}
O Lord, our Lord, how wonderful is Thy name over the whole earth! ℣. For thy
majesty is above the heavens.
\end{namedtranslation}

Here the chant and the Vulgate agree exactly. The whole Gradual is one
biblical verse, and the division into respond and verse falls on the
\latin{quóniam}, so that the responsorial form itself carries a logical
relation: the wonder, then its reason.

\begin{witnessnote}
\wlocus{St.\ Augustine, \work{Enarrationes in Psalmos} 8, §4}
Augustine treats the two halves as question and answer --- \textit{how} is
the Name wonderful throughout the earth? because the \latin{magnificentia} has
been lifted above the heavens --- and reads that lifting as the Ascension:
the glory was raised ``from earthly humiliation'' above the heavens, and it
was the visible ascent that disclosed who had descended. The preceding
sections have already read the psalm's title \latin{pro torcularibus},
``for the winepresses,'' of the Churches and the martyrdoms.
\end{witnessnote}

\begin{witnessnote}
\wlocus{St.\ Thomas Aquinas, \work{Super Psalmo} 8, n.~1}
Aquinas gives the sharpest account of why the second colon is a reason at all.
Wonder, he says, arises when the effect is seen and the cause unknown ---
\latin{Admiratio est quando aliquis videt effectum, et ignorat causam}. God is
not wholly unknown, since the creation shows him; but the effect does not
manifest the cause perfectly, and therefore he remains \latin{admirabilis}.
The exaltation above the heavens then excludes two errors at once: that God is
the form of the heavens, and that he acts by necessity of nature. He also
insists on \latin{in universa terra} against any restriction of the Name to
one territory.
\end{witnessnote}

\begin{witnessnote}
\wlocus{St.\ Robert Bellarmine, \work{Explanatio in Psalmos}, Ps.~VIII, v.~1}
Bellarmine anticipates the obvious objection --- that in fact few admire ---
and answers that the Name is \textit{worthy} of admiration whether or not
admirers are found, ``in the sense that all beautiful productions are said to
praise the producer.'' He then develops the second colon through an analogy of
princely magnificence: God ``created the universe for a palace, having the
earth for its pavement, the heavens for its roof.''
\end{witnessnote}

\textbf{Why this psalm answers this Epistle.} Psalm 8 is the psalm of human
smallness under a sky, and its continuation --- \latin{Quid est homo, quod
memor es ejus?} --- is quoted verbatim in Hebrews 2:6--8, where the argument
turns on the fact that we do not yet see all things subject to him. The
Gradual quotes only v.~2 and this guide claims no more than v.~2 for the
liturgy. But the appointed verse comes from a poem whose next move is exactly
the move the Epistle has just demanded: stop estimating your own standing.

\subsection*{Alleluia \cue{All.}}

\begin{appointedlatin}{Alleluia --- Ps. 58, 2 --- marginal no.\ 1526}
Allelúia, allelúia. ℣. \latin{Ps. 58, 2} Eripe me de inimícis meis, Deus meus:
et ab insurgéntibus in me líbera me. Allelúia.
\end{appointedlatin}

\begin{namedtranslation}{English: Shepherd, p.~229}
Alleluia, alleluia. ℣. Rescue me, O my God, from mine enemies: and, from them
that rise up against me, deliver me. Alleluia.
\end{namedtranslation}

The appointed verse matches the Clementine psalter word for word --- the only
one of the four psalm chants of which that is true without qualification.
Augustine, however, was reading an Old Latin text: his lemma runs \latin{Erue
me de inimicis meis, Deus meus, et ab insurgentibus super me redime me}, with
different verbs in both halves. His exegesis is of that text.

\textbf{An independent convergence worth reporting.} The psalm's title
refers the prayer to the night when Saul set a watch on David's house to kill
him. Three witnesses, in two languages and across eleven centuries, take that
guarded house to prefigure the guarded tomb.

\begin{receptiontable}{Witness}{Locus}{The reading}
St.\ Augustine & \work{Enarr.\ in Ps.} 58, Sermo I, §§1--4 & The title's \latin{ne disperdas \dots\ in tituli inscriptionem} is read of Pilate's inscription (\latin{Quod scripsi, scripsi}); Saul's watch on the house becomes the guard at the sepulchre, and Augustine presses the forensic absurdity of sleeping witnesses. §4 then applies the verse to the whole Christ: it was done once in his flesh and is done daily in his members\\
Theodoret of Cyrus & \work{Interpretatio in Psalmos} 58 & Narrates the title from 1 Sam.\ 18--19 in detail, then says the psalm refers only slightly to David's own time and forecasts the Passion: the guard set over the Lord's tomb found the garments but not the body\\
St.\ Robert Bellarmine & \work{Explanatio}, Ps.~LVIII, argument and v.~2 & Gives three applications in one sentence: David hemmed in by Saul's soldiers ``as if he were in a prison''; Christ ``as he lay in the sepulchre, with guards on it''; and ``any just person in danger of death''\\
\end{receptiontable}

Two things follow. The reading is not one author's idiosyncrasy: a Greek
Antiochene and a Latin African reached it independently, and a
sixteenth-century Doctor still had it. And it makes the Alleluia's placement
legible: the assembly sings, immediately before the Gospel of the approach to
Jerusalem, a prayer that the tradition hears in the mouth of the one who will
shortly be sealed in a watched tomb.

\textit{Negative result.} Aquinas wrote no commentary on this psalm. His
\work{Super Psalmos} breaks off at Psalm 54 with \latin{Deo gratias}, and
the appointed Alleluia therefore has no Thomistic exposition to compare.
Bellarmine is the later Doctoral witness in its place.

\subsection*{Gospel \cue{Gosp.}}

\begin{appointedlatin}{\latin{✠ Sequéntia sancti Evangélii secúndum Lucam} --- Luc. 19, 41-47 --- marginal no.\ 1527}
In illo témpore: Cum appropinquáret Iesus Ierúsalem, videns civitátem, flevit
super illam, dicens: Quia si cognovísses et tu, et quidem in hac die tua, quæ
ad pacem tibi, nunc autem abscóndita sunt ab óculis tuis. Quia vénient dies in
te: et circúmdabunt te inimíci tui vallo, et circúmdabunt te: et coangustábunt
te úndique: et ad terram prostérnent te, et fílios tuos, qui in te sunt, et non
relínquent in te lápidem super lápidem: eo quod non cognóveris tempus
visitatiónis tuæ. Et ingréssus in templum, cœpit eícere vendéntes in illo, et
eméntes, dicens illis: Scriptum est: Quia domus mea domus oratiónis est. Vos
autem fecístis illam spelúncam latrónum. Et erat docens cotídie in templo.
\end{appointedlatin}

\begin{namedtranslation}{English: Shepherd, pp.~229--230}
At that time: When he drew near Jerusalem, seeing the city, he wept over it,
saying: If thou also hadst known, and that in this thy day, the things that are
to thy peace: but now they are hidden from thy eyes. For the days shall come
upon thee: and thy enemies shall cast a trench about thee, and compass thee
round, and straiten thee on every side, and beat thee flat to the ground, and
thy children who are in thee: and they shall not leave in thee a stone upon a
stone: because thou hast not known the time of thy visitation. And entering into
the temple, he began to cast out them that sold therein, and them that bought.
Saying to them: It is written: My house is the house of prayer: but you have made
it a den of thieves. And he was teaching daily in the temple.
\end{namedtranslation}

\subsubsection*{The two prophets fused in one sentence}

Christ's Temple saying joins two texts that pull in different directions, and
both contexts matter.

\begin{receptiontable}{Text}{Vulgate wording}{Its own literary setting}
Isaiah 56:7 & \latin{quia domus mea domus orationis vocabitur cunctis populis} & The climax of an oracle of \textit{inclusion}: the eunuch is promised a name better than sons and daughters, and the foreigner who holds fast to the Lord is promised that his sacrifices will be accepted on the altar. What follows immediately (56:9--12) is a savage indictment of Israel's own leaders --- blind watchmen, dumb dogs, each turned to his own gain\\
Jeremiah 7:11 & \latin{Numquid ergo spelunca latronum facta est domus ista, in qua invocatum est nomen meum} & The Temple Sermon, preached in the gate of the Lord's house against the belief that the sanctuary confers immunity: \latin{Nolite confidere in verbis mendacii, dicentes: Templum Domini, templum Domini}. The charge sheet is theft, murder, adultery, perjury and Baal; the threat is that God will do to this house as he did to Shiloh\\
\end{receptiontable}

Set together, the two quotations say something precise: the house whose
vocation is to gather all peoples has been made into the place where the
guilty feel safe. And Luke quotes the first of them \textit{short}. Mark
11:17 keeps \latin{omnibus gentibus}; Luke does not. The Roman Office of
Guéranger's day restored the fuller form in its Magnificat antiphon for this
Sunday --- \latin{domus orationis est cunctis gentibus} --- which is a
liturgical conflation of Luke with Mark and Isaiah and must not be cited as
Luke's text.

\subsubsection*{The tears}

\begin{witnessnote}
\wlocus{Origen, \work{Homiliae in Lucam} 38, in Jerome's Latin (PL 26, cols.\ 302B--303D)}
Origen begins from a principle: Christ fulfils in himself every beatitude he
preached, and therefore he wept in order to found the beatitude of those who
mourn. He then stages an objector who says the prophecy was plainly fulfilled
when the Roman army destroyed the city --- concedes it --- and refuses to
stop there. \latin{Nos enim sumus Jerusalem quae defletur}: the weeping
concerns \textit{our} Jerusalem, the baptized who sins after the mysteries.
``He does not lament over a Gentile, but over him who was a citizen of
Jerusalem and ceased to be so.'' Consequently the besiegers are
demonologised: \latin{circumdant eam inimici, contrariae videlicet
fortitudines, spiritus nequam} --- especially, he adds, when someone is
overcome after years of chastity. \textit{Rights note:} the only English
translation of this homily is modern and in copyright; the Latin is
public domain, so the Latin is quoted and the rest is analysis.
\end{witnessnote}

\begin{witnessnote}
\wlocus{St.\ Cyril of Alexandria, \work{Commentary on Luke}, Sermon 131, tr.\ R.~Payne Smith (Oxford, 1859)}
Cyril's interest is what a tear proves. The ruin was not God's good pleasure;
Christ pitied them. The weeping is a real human act performed as the
disclosure of an otherwise invisible compassion --- ``that we hereby might
learn that He feels grief, if we may so speak of God, Who transcends all''
--- and ``the tear which drops from the eye is a symbol of grief, or rather, a
plain demonstration of it.'' \latin{Quæ ad pacem tibi} he glosses as ``the
things, that is, useful and necessary for you to make your peace with God,''
and he reads the hiding through 2 Corinthians 3 (the veil) and Romans 11:25,
where the blindness is \textit{in part}. He treats the siege as fulfilled ---
without ever naming Titus, Vespasian or the Romans --- and immediately
softens it: Israel ``did not perish from the very roots,'' the apostles being
the firstfruits.
\end{witnessnote}

\begin{witnessnote}
\wlocus{St.\ Gregory the Great, \work{Homiliae in Evangelia} 39, §§1--2 (PL 76, 1293--1301)}
Gregory preached this homily to the people in the Lateran basilica on this
very pericope, and his opening sentence is the one the later tradition
repeats: \latin{Flevit etenim prius Redemptor ruinam perfidae civitatis, quam
ipsa sibi civitas non cognoscebat esse venturam} --- ``The merciful Redeemer
wept over the fall of that city, which the city itself did not know was
coming.'' He supplies the elided verb of the broken sentence: \latin{si
cognovisses \dots\ subaudi, fleres} --- understand, ``thou wouldst have
wept.'' He names the agents of the destruction, Vespasian and Titus, and
argues that not one stone was left from the city's own later migration to the
ground outside the gate where the Lord was crucified. The cause is the
unrecognised visitation, clinched with Jeremiah 8:7, where kite, turtledove,
swallow and stork all keep their seasons.
\end{witnessnote}

\subsubsection*{The second pass: Gregory turns the whole pericope inward}

At §3 Gregory announces a deliberate change of register --- \latin{debemus ex
rebus exterioribus introrsus aliquam similitudinem trahere} --- and then
reruns the entire text of the soul. Christ wept once over the city, but
``does not cease to do this daily through his elect'' over those who fall from
a good life. The soul's \latin{dies sua} is its transitory time. The
besieging enemies are the demons --- \latin{Qui unquam sunt humanae animae
majores inimici, quam maligni spiritus?} --- who encircle the dying soul with
a rampart built out of its own remembered iniquities; the children dashed to
the ground are its illicit thoughts; stone upon stone is perverse thought
heaped on perverse thought. At §5 he says God visits every soul in three
ways: \latin{praecepto, flagello, miraculo}.

This second pass is where the homily's weight falls, and it is worth noticing
that the Roman Breviary's Matins lessons for this Sunday stop at §§1--2 ---
the historical reading --- and never reach it.

\subsubsection*{The Temple, and who the tradition thinks the thieves are}

The striking fact about the patristic handling of \latin{spelúncam latrónum}
is how consistently it is turned on the clergy.

\begin{receptiontable}{Witness and locus}{The reading of the Temple action}{Whom it is aimed at}
St.\ Ambrose, \work{Expositio Ev.\ sec.\ Lucam} IX.16--22 (PL 15) & \latin{Deus templum suum non mercatoris vult esse diversorium, sed domicilium sanctitatis}. The money-changers are those who traffic in Scripture --- ``the Lord's money is divine Scripture'' --- and not all changers are expelled, since Mt.\ 25:27 knows good ones & Preachers who debase the King's image; and, at §19, those who cannot belong in the Church because they \latin{sancti Spiritus gratiam nundinentur} --- who put the grace of the Holy Spirit up for sale\\
Origen, \work{Hom.\ in Lucam} 38 & Presses the letter: Jesus expels only \textit{sellers}, not buyers, and the only creature named is the dove & The dove is the Holy Spirit; therefore to sell doves is to teach the revealed mysteries for a fee\\
St.\ Gregory, \work{Hom.\ in Ev.} 39, §§2, 6--7 & \latin{ruina populi maxime ex culpa sacerdotum fuit} --- the ruin of the people came chiefly from the fault of the priests. The \latin{spelunca latronum} is bribery in the Temple: those who did not pay were persecuted in the body, those who did were killed in the spirit & Explicitly those ``who confer the imposition of hands for a fee''; then, at §6, the \latin{vita religiosorum} that sells justice, and at §7 the conscience of any believer that plots harm to a neighbour\\
\end{receptiontable}

Three cautions belong with that table, and each corrects a claim it would be
easy to make.

\textbf{Ambrose is not commenting on Luke here.} His lemma at §18 is Matthew
21:12 and at §§20--21 John 2:15; he never expounds Luke's wording. More
important, Ambrose has no exegesis of Luke 19:41--44 anywhere: neither in the
\work{Expositio}, which goes straight from the triumphal entry at §16 to the
Temple at §17, nor elsewhere in his corpus. He also never quotes \latin{domus
orationis} or \latin{spelunca latronum}.

\textbf{Ambrose also says something the tradition must not be allowed to
hide.} In the same section, at §22, he writes that the Jews are expelled
\latin{non praescripto aliquo exsilio, sed intermino; ut nusquam Synagogae
locus in orbe remaneret}. That is a supersessionist sentence, it is in this
very passage, and suppressing it would falsify the reception history this
guide is reporting.

\textbf{Bede's exposition of this pericope is largely Gregory's.} Bede,
\work{In Lucae Evangelium Expositio} V (PL 92, 570C--574B), reproduces
Gregory's Homily 39 closely enough that whole clauses recur verbatim ---
\latin{ruinam perfidae civitatis}, \latin{ex culpa sacerdotum fuit},
\latin{Sicut templum Dei in civitate est}, \latin{quotidie Veritas in templo
docet}. Bede told Acca that he marked patristic borrowings with marginal
initials; Migne does not preserve them. What \textit{is} Bede's own is the
closing block at 573C--574B: an Exodus 12 typology in which the lamb taken on
the tenth day corresponds to the five days between the reception of the Lord
near the Mount of Olives and the Pasch --- lamb because he takes away sin, kid
because he was charged with it; the crowd singing \latin{Benedictus qui venit}
brought in the lamb, and the Pharisees' \latin{Magister, increpa discipulos
tuos} the kid. This is the pericope's own bridge to the Passion, and it is
Bede's.

\subsubsection*{Two large negative results}

\textbf{St.\ Augustine has no exegesis of Luke 19:41--44.} A search of the
sermon corpus, all 150 \work{Enarrationes}, the 124 \work{Tractates on John},
all twenty-two books of \work{De civitate Dei}, \work{De consensu
evangelistarum}, the \work{Quaestiones Evangeliorum} and the letters produced
two loci only: \work{Epistula} 199 (\work{De fine saeculi}), cap.~IV §12,
where he is \textit{replying} to Hesychius's use of the verse and insists that
it concerns the first advent and not the second --- \latin{sed hoc ad tempus
primi aduentus domini pertinet, non secundi} --- and one clause in \work{De
sancta virginitate} 28: ``Blessed are they that mourn; imitate Him, Who wept
over Jerusalem.'' \work{De consensu evangelistarum} II.66 harmonises the
triumphal entry and moves straight to the cleansing at II.67, passing over
Luke's unique pericope in the one work devoted to Lucan singularities; and the
\work{Quaestiones Evangeliorum} II runs qq.~XLVI--XLVIII, then an unnumbered
question on Luke 19:45--46, then q.~XLIX, with no question on 19:41--44 at all.
Augustine \textit{does} treat \latin{domus orationis} against \latin{spelunca
latronum} at \work{Enarr.\ in Ps.} 130 (LXX) §§1--3, and the sentence there is
the one this guide would most like to have had from him on the Gospel:
\latin{quantum in ipsis est, domum Dei volunt facere speluncam latronum; nec
ideo evertunt templum} --- so far as in them lies they want to make God's
house a robbers' den, and they do not thereby overthrow the temple.

\textbf{St.\ Thomas never cites this Gospel in the \work{Summa}.}
\latin{Luc.\ XIX} occurs four times in the whole work, each time a different
verse (19:8, 19:23, the pounds, 19:10); \latin{flevit} occurs once, inside a
quotation from Hilary at III q.15 a.5 ad~1; \latin{spelunca latronum} and
\latin{videns civitatem} occur not at all. Where he might have used the tears
--- III q.15 a.6, on sorrow in Christ --- his \latin{sed contra} is Matthew
26:38. Where he might have used the cleansing --- III q.15 a.9, on anger in
Christ --- his \latin{sed contra} is John 2:17. Where he treats simony,
II--II q.100, the temple-cleansing texts do not appear at all. Thomas's
engagement with this pericope is real but indirect, and it runs through the
\work{Catena aurea}, which is a compilation and not his exegesis.

\subsubsection*{What the \work{Catena aurea} shows about the shape of the tradition}

The \work{Catena} divides the pericope at v.~44 and assembles thirteen
excerpts for vv.~41--44 and fifteen for vv.~45--48. Gregory supplies six of
the first thirteen and five of the second fifteen, and in both lections he
supplies the closing tropology; the Greeks --- Origen, Cyril, Eusebius,
Theophylact --- carry the literal and historical sense. Ambrose, Bede and
Augustine appear only in the Temple half, which independently corroborates
the three findings above. The compilation is a map of reception, not a
witness: the 1843 Oxford edition's preface warns that for St.~Luke the Greek
citations reached Aquinas through translations made for him which often gave
``not always the very words \dots\ but frequently only the sense.''

\subsubsection*{The reception's real danger, and the Church's own correction}

The tradition summarised above contains a strain that has done grave harm,
and a guide that reproduced the strain silently would be dishonest twice
over. The most convenient witness is the same book this guide has been
quoting for its English. Shepherd's 1909 volume introduces the Mass of this
Sunday with ``Israel had made himself the enemy of the Church; and God \dots\
punishes and disperses his children'' (p.~224), and on the Gospel writes of
``the Jewish deicides'' and ``the apostate nation'' (pp.~227, 230). Those are
the words of the translation whose English propers appear in every block
above. They are printed here so the reader knows what the source is.

That reading has since been named and rejected by the Church's own organs, and
the correction bears directly on this pericope.

\begin{receptiontable}{Document and locus}{What it says}{Bearing on this Gospel}
\work{Nostra aetate} 4 (1965) & Writes that ``Jerusalem did not recognize the time of her visitation,'' attaching note~9, \latin{Cf. Lk.\ 19:44}; then states that ``the Jews should not be presented as rejected or accursed by God, as if this followed from the Holy Scriptures,'' and decries ``hatred, persecutions, displays of anti-Semitism, directed against Jews at any time and by anyone'' & The Council cites the appointed verse and blocks the inference inside the same article. That adjacency is the model this guide follows\\
\work{Notes on the Correct Way to Present the Jews and Judaism} (1985), VI.1 & ``The history of Israel did not end in 70 A.D.'' and ``We must rid ourselves of the traditional idea of a people \textit{punished}, preserved as a \textit{living argument} for Christian apologetic'' & Names precisely the reading Shepherd's commentary embodies, and rejects it\\
Pontifical Biblical Commission, \work{The Jewish People and Their Sacred Scriptures in the Christian Bible} (2001/2002), §§51, 53, 74 & §51: he ``tearfully foresees'' the ruin. §53: ``The divine sanction will be the same as in Jeremiah's time,'' followed at once by ``as in Jeremiah's time --- God is not satisfied merely to punish, he also offers pardon.'' §74: Luke shows Jesus weeping over Jerusalem (19:41--44) and disregarding his own sufferings for the women and children of that city (23:28--31); ``universalism does not mean being anti-Jewish'' & Supplies the controlling analogy: Jeremiah. The pattern is intra-covenantal prophetic warning, which is exactly what Jeremiah 7 --- the text Christ quotes --- already was\\
\end{receptiontable}

The point is not that the Fathers must be edited. It is that the pericope's own
quotations settle the register. Christ indicts the Temple with Jeremiah's
Temple Sermon, and Jeremiah was a prophet of Israel speaking inside Israel
about Israel's own sanctuary. The Epistle read ten minutes earlier has already
said whose correction the whole thing is: \latin{ad correptiónem nostram}.

\subsubsection*{The last clause}

\latin{Et erat docens cotídie in templo.} The imperfect periphrastic is
durative: not ``he taught'' but ``he was teaching, daily.'' Gregory reads it
as grace following discipline, and glosses it in the moral pass as
\latin{quotidie Veritas in templo docet} --- when Truth patiently instructs
the mind of the faithful to avoid evil, Truth is teaching daily in the temple.
The missal's decision to end the reading here, rather than four words later
where the plot begins, means the assembly's last impression of the day's
Gospel is of a teacher who did not leave.

\subsection*{Offertory \cue{Off.}}

\begin{appointedlatin}{\latin{Antiphona ad Offertorium} --- Ps. 18, 9, 10, 11 et 12 --- marginal no.\ 1528}
Iustítiæ Dómini rectæ, lætificántes corda, et iudícia eius dulcióra super mel
et favum: nam et servus tuus custódit ea.
\end{appointedlatin}

\begin{namedtranslation}{English: Shepherd, p.~249}
The justices of the Lord are right, rejoicing hearts; and his precepts are
sweeter than honey and the honey-comb; and therefore doth thy servant observe
them.
\end{namedtranslation}

The antiphon is built, not quoted. It takes v.~9a, carries over the subject
\latin{iudicia} from v.~10c, borrows the predicate \latin{dulciora super mel
et favum} from v.~11b, and closes with v.~12a, replacing the psalm's
\latin{Etenim} with \latin{nam et}. That is why the missal prints four verse
numbers instead of a range. The construction is faithful to the psalm's own
grammar --- in vv.~10--11 it is the judgments that are more desirable than
gold and sweeter than honey --- but the appointed text is a made thing, and
this guide says so rather than describing it as ``Psalm 18:9--12.''

\begin{witnessnote}
\wlocus{St.\ Thomas Aquinas, \work{Super Psalmo} 18, nn.~5--7}
Aquinas predicates two things of the \latin{iustitiae}: they are \latin{recta},
containing justice, and \latin{iucunda}, ``not severe and disturbing, because
they are mixed with equity'' --- hence \latin{laetificantes corda}, on account
of that equity and the hope of reward. At n.~7 he reaches the honey and gives
three reasons why spiritual delights are objectively greater than bodily ones:
the good enjoyed is higher; the power enjoying (intellect) is greater than
sense; and the mode differs, since bodily pleasures consist \latin{in fieri et
in motu}, and what is only in becoming is never possessed whole at once. The
proof, he says, is twofold --- by experience, \latin{nam et servus tuus
custódit ea}, and by effect. Then he presses the preposition: Scripture does
not say \latin{pro custodia} but \latin{in custodiendis illis}, because the
keeping of them is itself the great reward.
\end{witnessnote}

\begin{witnessnote}
\wlocus{St.\ Augustine, \work{Enarr.\ in Ps.} 18, Sermo I §§9--12 and Sermo II §§9--12}
Sermo I reads the law-section christologically --- the spotless law is Christ
--- and produces the image the Offertory most invites: one may already be
honey, freed from this life, or still be honeycomb, wrapped in this life as in
wax, needing God's hand not to crush but to press one out of temporal into
eternal life. At §12 he insists the reward for keeping the judgments is not
some external benefit but the keeping itself. Sermo~II reads the same verses
pneumatologically and against the Donatists, punctuating each clause with
\latin{Hoc est Spiritus sanctus} and turning \latin{iustificata in idipsum}
onto ecclesial unity: \latin{non ad rixas divisionis, sed ad congregationem
unitatis}. \textit{Access limit:} the second sermon is not in the
public-domain NPNF English, which prints only the first; it was read in Latin.
\end{witnessnote}

\begin{witnessnote}
\wlocus{Theodoret of Cyrus, \work{Interpretatio in Psalmos} 18 (PG 80, cols.\ 996--997)}
Theodoret expounds the five names the Law gives itself --- law, testimony,
judgments, command, decrees --- and assigns each a distinct function, the
judgments gladdening the heart by revealing the basis of judgment. He then
qualifies the sweetness in a way no other witness here does: these things are
worth more than gold and sweeter than honey not to all human beings, but to
those truly human, whose life is not comparable with the beasts. Summarised,
not quoted, for the rights reason given above.
\end{witnessnote}

\textbf{A variant that changes the mood.} Dom Dominic Johner, comparing the
Missal with Dom Pothier's \work{Liber Gradualis} of 1883 and 1895, observes
that the older chant books close this Offertory in the sixth mode and require
\latin{custodiet} --- ``he \textit{will} keep them,'' a resolve --- whereas the
Missal prints the indicative \latin{custódit}. On the older reading, he says,
the melody's close on the half tone ``is far removed from victorious
certainty. It sounds like a fervent petition.'' Of this Sunday specifically he
writes only that, after the threatened destruction of the city in the Gospel,
the piece is to be sung ``in a somewhat more subdued fashion.''
\textit{Rights note:} Johner's book was published in the United States in
1940; a search of the standard renewal record found no renewal, which points
to but does not establish public-domain status, so he is paraphrased and
quoted only in short compass.

\subsection*{Secret \cue{Sec.}}

\begin{appointedlatin}{\latin{Secreta} --- marginal no.\ 1529}
Concéde nobis, quæsumus, Dómine, hæc digne frequentáre mystéria: quia,
quóties huius hóstiæ commemorátio celebrátur, opus nostræ redemptiónis
exercétur. Per Dóminum nostrum Iesum Christum, Fílium tuum: Qui tecum vivit et
regnat in unitáte.
\end{appointedlatin}

\begin{namedtranslation}{English: Shepherd, p.~249}
Grant us, O Lord, we beseech thee, frequently and worthily to celebrate these
mysteries: for, as many times as this commemorative sacrifice is celebrated, so
often is the work of our redemption performed. Through, \textit{etc}.
\end{namedtranslation}

The prayer's whole weight is on a \latin{quia}. It does not ask for worthy
celebration and then add a pious reason; the reason is the premise, and the
petition follows from it. Three words carry it. \latin{Frequentáre} does not
mean ``to frequent'' in the modern sense but to do a thing repeatedly and
attentively --- Shepherd's ``frequently and worthily to celebrate'' splits it
into two English words. \latin{Quóties} makes the relation distributive: not
``because this is a great mystery'' but ``as often as.'' And
\latin{exercétur} is present passive: the work is being carried on, now,
here, in the celebration.

\textbf{This is the one appointed text of the day that the Church has quoted
back.} \work{Sacrosanctum Concilium} 2 writes that in the liturgy, and
supremely in the Eucharistic Sacrifice, \latin{opus nostrae Redemptionis
exercetur}, and its note~1 reads \latin{Missale romanum, oratio super oblata
dominicae IX post Pentecosten}. \work{Presbyterorum ordinis} 13 reuses the
clause --- ``In mysterio Sacrificii Eucharistici, in quo munus suum praecipuum
sacerdotes adimplent, opus nostrae redemptionis continuo exercetur'' --- and
its note 105 identifies the same prayer. The two conciliar uses are not
identical: the Decree inserts \latin{continuo}, so it is documented reuse
rather than verbatim repetition, and this guide reports it as such.

The typical edition prints the long conclusion here --- \latin{Per Dóminum
nostrum Iesum Christum, Fílium tuum: Qui tecum vivit et regnat in unitáte} ---
where Benziger prints only \latin{Per Dóminum}. Since the Secret is said
silently and this conclusion is the point at which the celebrant's voice
returns aloud, the difference is not merely typographical for anyone using the
book.

\subsection*{Communion \cue{Comm.}}

\begin{appointedlatin}{\latin{Ant. ad Communionem} --- Io. 6, 57 --- marginal no.\ 1530}
Qui mandúcat meam carnem, et bibit meum sánguinem, in me manet, et ego in eo,
dicit Dóminus.
\end{appointedlatin}

\begin{namedtranslation}{English: Shepherd, p.~250}
He that eateth my flesh, and drinketh my blood, abideth in me, and I in him,
saith the Lord.
\end{namedtranslation}

\textbf{Two divergences from the Bible on the page.} The antiphon adds three
words to the Johannine verse --- \latin{dicit Dóminus} --- which are the
liturgy's own attribution formula, added so that a first-person sentence can
be sung by a choir without the singers claiming it. And it reads \latin{et ego
in eo} where the Clementine Vulgate reads \latin{et ego in illo}. Both were
checked directly against the Clementine text of John 6:57.

\textbf{Why the verse number differs between editions, demonstrated rather
than asserted.} The Clementine chapter has seventy-two verses; the critical
text has seventy-one. The offset begins at critical v.~51, which the Clementine
divides in two: its v.~51 is \latin{Ego sum panis vivus, qui de caelo
descendi} and its v.~52 begins \latin{Si quis manducaverit ex hoc pane}. From
that point on, Clementine equals critical plus one --- which is exactly why the
missal prints \latin{Io. 6, 57} for a clause most current editions number
6:56. The guide keeps the missal's number and states the other.

\textbf{The grammar.} The clause is a chiasm of persons. \latin{In me manet}
moves the communicant into Christ; \latin{et ego in eo} moves Christ into the
communicant; the verb is not repeated, so the second half hangs grammatically
on the first. One act of eating, two directions of indwelling. Of the ten
propers this is the only one whose motion is reciprocal: every other text of
the day runs either from God toward the assembly or from the assembly toward
God.

\begin{witnessnote}
\wlocus{St.\ Augustine, \work{Tractates on John} 26, §§11--18, and 27, §§1, 6}
Augustine builds the apparatus before he touches the clause. At §12 he
separates the visible rite from its power --- ``the sacrament is one thing, the
virtue of the sacrament another'' --- and gives the formula the whole later
tradition uses: ``he that eats within, not without; who eats in his heart, not
who presses with his teeth.'' At §13 he brings in 1 Corinthians 10:17, the
verse the day's own Epistle stops four verses short of: ``One bread, says he,
we being many are one body. O mystery of piety! O sign of unity! O bond of
charity!'' At §15 he distinguishes the sacrament of the thing from the thing
itself, which is the fellowship of the body and its members. Only then, at
§18, does he reach the antiphon's clause: ``This it is, therefore, for a man to
eat that meat and to drink that drink, to dwell in Christ, and to have Christ
dwelling in him.'' Tractate 27 §1 turns it into a test: ``The proof that a man
has eaten and drank is this, if he abides and is abode in.'' English from the
public-domain NPNF first series, vol.~7, tr.\ John Gibb. \textit{Correction of
a common citation:} 1 Corinthians 10:17 is quoted in Tractate 26 only, not in
Tractate 27.
\end{witnessnote}

\begin{witnessnote}
\wlocus{St.\ John Chrysostom, \work{Homilies on John} 47}
Chrysostom's word for the union is mixture: ``He that eats My flesh, dwells in
Me. This He said, showing that such an one is blended with Him.'' The
consequence follows as a syllogism --- ``for if he dwells in Me, and I live, it
is plain that he will live also.'' He then guards the following verse against
a subordinationist reading and insists that the life promised is not bare life,
since unbelievers live too, but ``that glorious and ineffable life.'' NPNF
first series, vol.~14, tr.\ Charles Marriott.
\end{witnessnote}

\begin{witnessnote}
\wlocus{St.\ Cyril of Alexandria, \work{Commentary on John} IV.2, tr.\ P.~E.~Pusey (Library of the Fathers, 1874)}
Cyril illustrates the indwelling twice, physically: ``as if one should join wax
with other wax, he will surely see (I suppose) the one in the other,'' and
then leaven in the lump, so that ``the least portion of the Blessing blendeth
our whole body with itself.'' He draws a pastoral conclusion that runs
\textit{against} abstention: those who plead unworthiness will, if they wait
until they have stopped stumbling, never communicate at all, since the
Blessing is medicinal.
\end{witnessnote}

\begin{witnessnote}
\wlocus{St.\ Thomas Aquinas, \work{Super Ioannem} c.~6, lect.~7, n.~976 --- and this one is authentic}
Thomas casts the passage as a syllogism in which the antiphon's verse is the
major premise: whoever eats is joined to Christ; whoever is joined to Christ
has life. Taken \latin{mystice} there is no difficulty --- spiritual eating is
incorporation \latin{per unionem fidei et caritatis}. Taken of sacramental
reception he distinguishes, on Augustine's authority (\work{Tract.\ in Ioh.}
26.18, quoted above), those who eat \latin{non sacramentaliter tantum, sed
revera} from those \latin{qui in corde ficto ad illud accedunt}, in whom
\latin{nullum enim effectum habet sacramentum in ficto}. And he defines the
feigning structurally: \latin{Fictus enim est, cum non respondet interius quod
signatur exterius} --- and what this sacrament outwardly signifies is precisely
the mutual indwelling. So the appointed clause functions in Thomas as the
criterion by which fruitful reception is told from bare reception.
\end{witnessnote}

Johner hears the same structure in the melody, and the observation is worth
recording because it is arrived at independently of the exegesis: the
descending fourth at \latin{in me manet} is answered by a rising fourth at
\latin{ego in eo}, ``thus both thoughts are placed in strong relief: Thou in me
and I in Thee.''

\subsection*{Postcommunion \cue{Postcomm.}}

\begin{appointedlatin}{\latin{Postcommunio} --- marginal no.\ 1531}
Tui nobis, quæsumus, Dómine, commúnio sacraménti, et purificatiónem cónferat,
et tríbuat unitátem. Per Dóminum nostrum.
\end{appointedlatin}

\begin{namedtranslation}{English: Shepherd, p.~250}
May the participation of this thy sacrament, O Lord, we beseech thee, both
purify us, and unite us. Through, \textit{etc}.
\end{namedtranslation}

The Latin is a single sentence with one subject and two verbs, and the
\latin{et \dots\ et} construction makes them equally weighted: not purification
\textit{and then} unity, but both at once from one act. The genitive phrase
\latin{Tui \dots\ sacraménti} is split around the address, which throws
\latin{Tui} into first position: the sacrament is \textit{thine} before it is
ours.

Two things about the pairing repay attention. First, the two objects match
the two fruits the Sunday has been moving toward: purification answers the
Epistle's warning and the Gospel's cleared Temple, and unity answers the
Communion's mutual abiding. Second, neither is described as accomplished. The
formulary that opened by asking God to re-form our desires closes by asking
for a cleansing and a unity that are still being asked for. Shepherd's own
comment on this prayer is exact: ``The sanctification of each individual
member of the Church, and the unity of the social body, are the two fruits of
these sacred mysteries.''
